\chapter{Introduction}
\label{ch:1}
\noindent
\section{Motivation}
Mass spectrometry (MS) is one of the most widely used analytical techniques for molecular characterization across chemistry, metabolomics, environmental sciences, and pharmaceutical research. \cite[p. 1]{Gross2006}
Its high sensitivity \cite[p.~1764]{Cramer2020}, reproducibility \cite[p.~1764]{Cramer2020}, and scalability \cite[p.~1764]{Cramer2020} 
have enabled large-scale molecular profiling in both research and industrial settings.
Modern instrumentation such as Orbitrap, Q-TOF, and triple quadrupole mass spectrometers 
generates high-throughput spectral data that can be compared against large repositories of molecular fingerprints 
such as the NIST Mass Spectral Library to support compound identification and profiling \cite[pp.~7274-7275]{Stein2012}.

Despite these advances in data acquisition, the interpretation of MS data remains a significant bottleneck \cite[p.~1]{Blazenovic2018}. 
Structural elucidation from spectra is still largely expert-driven, relying on chemical intuition and established fragmentation rules \cite[p.~11]{Blazenovic2018}. 
While thousands of spectra can be recorded within a short time frame \cite[p.~1]{Blazenovic2018}, the corresponding structural annotation process remains comparatively slow, manual, and difficult to scale \cite[p.~3]{Blazenovic2018}. 
This imbalance creates a fundamental asymmetry between data generation and data analysis.

The intrinsic complexity of molecular fragmentation contributes substantially to this challenge \cite[p.~240]{Gross2006}. 
Fragmentation processes are governed by chemically valid bond cleavages and rearrangements \cite[pp.~30-31]{Gross2006},
%which are determined by internal energy redistribution and bond dissociation energetics \cite[p.~24-26]{Gross2006}, 
leading to multiple competing fragmentation pathways. 
A single molecule may therefore generate numerous fragments through direct cleavages and rearrangement reactions \cite[p.~240-241]{Gross2006}. 
Because fragmentation pathways depend sensitively on molecular structure and internal energy \cite[p.~26,pp.~30-31]{Gross2006}, 
structurally distinct molecules may produce partially overlapping spectral signatures, while small structural variations can induce significant spectral differences. 
Consequently, the inverse problem --- inferring molecular structure from spectral data --- is inherently ill-posed \cite[p.~240]{Gross2006}.

Recent machine learning approaches aim to address this challenge by learning direct mappings between molecular representations and spectral outputs. 
Graph neural networks \cite{Young2024}, sequence-based encoders \cite[p.~1573]{Schwaller2019}, and generative models have demonstrated promising performance in forward (molecule-to-spectrum) and inverse (spectrum-to-molecule) prediction tasks. 
However, purely data-driven approaches exhibit several limitations:
\begin{itemize}
    \item They often lack mechanistic interpretability. \cite[p.~3\&~6]{Nowatzky2025}
    \item They do not explicitly model fragmentation pathways. \cite[p.~3]{Nowatzky2025}
    \item They rely on statistical correlations rather than chemically structured priors. \cite[p.~966]{Goldman2023}
    \item Their generalization to rare or out-of-distribution molecular scaffolds remains limited. \cite[p.~12]{Russo2024}
\end{itemize}

In particular, models that treat molecules as flat graphs or string encodings and spectra as unordered peak lists ignore the structured chemical processes underlying fragmentation. 
This omission reduces interpretability and may hinder robust extrapolation beyond the training distribution.
\cite[p.~3420]{Goldman2024}

Fragmentation in mass spectrometry is not an arbitrary stochastic process; 
it follows chemically constrained transformation sequences. \cite[ch.~4\&~8]{McLafferty1993}
Therefore, incorporating structured representations of fragmentation—such as derivation graphs encoding 
valid transformation rules—offers a principled way to introduce chemical inductive bias into machine learning models. 
By explicitly modeling fragment lineage and transformation pathways, 
it becomes possible to align molecular substructures with spectral evidence in a mechanistically meaningful manner.


This thesis is motivated by the hypothesis that integrating chemically structured derivation graphs into a machine learning framework can improve predictive robustness, interpretability, and bidirectional structure-spectrum modeling. By combining fragment-aware representations with data-driven learning, the aim is to bridge the gap between high-throughput spectral acquisition and scalable, chemically grounded interpretation.

\section{Objectives}
The objective of this thesis is to design, implement, 
and systematically evaluate a fragment-augmented machine learning framework 
for improving the bidirectional translation 
between molecular structures and mass spectra. 
Specifically, the work investigates how the integration of graph-based reaction 
and fragment derivation graphs, generated through chemically valid graph transformation rules, 
can serve as structured and interpretable intermediate representations 
that enhance both predictive performance and model transparency.

Mass spectrometry-based structure elucidation represents an inherently ill-posed inverse problem: 
structurally distinct molecules may yield partially overlapping spectral signatures, 
while minor structural modifications can induce substantial spectral changes. 
This thesis therefore introduces fragment-level intermediates derived from rule-based graph transformations 
to inject chemically grounded structure into the learning process.

Instead of directly learning a mapping between molecular graphs and spectral representations, 
the proposed framework decomposes the task into a structured pipeline in which molecular graphs 
are transformed into chemically plausible fragment sets, which in turn mediate the relationship to spectral outputs.
These fragment sets are embedded into a shared latent space together with molecular and spectral representations, 
enabling aligned embeddings across the domains.

Ultimately, this thesis aims to demonstrate that incorporating chemically 
valid graph-based fragment derivations into neural learning pipelines, 
enhancing both the predictive power and the interpretability 
of bidirectional structure-spectrum modeling in mass spectrometry 
and enabling de novo molecular generation from spectral data.


\section{Scope and Limitations}

This thesis is restricted to electron ionization (EI) mass spectrometry, 
as EI fragmentation reactions are comparatively 
well understood and mechanistically characterized. \cite[ch.~4\&~8]{McLafferty1993}
This allows the use of chemically informed rule-based graph transformations 
to generate fragment derivation graphs under consistent assumptions. 
Other ionization techniques are excluded due to 
their differing fragmentation behavior 
and the additional modeling complexity they would introduce.

A major limitation arises from the computational cost of graph-based fragment generation. 
Rule-based graph transformations lead to combinatorial growth in possible fragments, \cite[p.~4]{Andersen2013}
especially when multiple reaction rules and recursive fragmentation steps are considered. 
Ensuring chemical validity and avoiding duplicate graph states 
further increases computational demands. \cite[pp.~2-3]{Boecker2016}

Due to this scalability constraint, the experimental evaluation is 
limited to order of magnitude of 100 compounds. 
While sufficient for demonstrating methodological feasibility and analyzing interpretability, this restricted dataset 
limits chemical diversity and does not allow for generalization. 
Consequently, the results should be interpreted as proof-of-concept validation 
rather than evidence of large-scale applicability.

It is also important to note that the reactions are defined in terms of graph symbols in an ambiguous manner.
Furthermore, the transformation rules are not exhaustive of all possible fragmentation pathways, and some may even be incorrect.
This means that the generated fragment derivation graphs may not capture all relevant chemical transformations, 
and the model's performance may be sensitive to the specific choice of transformation rules. 
Therefore, the interpretability analysis is constrained by the assumptions embedded in the rule set.
Consequently, it may not generalize to fragmentation processes that deviate from these assumptions.